% TOP_PROBLEM: {"id": 30007044, "title": "Maximal multipoint Seshadri constants in higher complex dimension", "category_id": 5, "category": {"id": 5, "name": "algebraic_geometry", "display_name": "Algebraic Geometry", "description": "Geometric objects defined by polynomial equations.", "slug": "algebraic-geometry", "order_index": 5, "created_at": "2026-07-31T15:26:25.671Z"}, "status": "open"}
% FIRST_POSTED: 2026-10-06
% ATTEMPT_STATUS: solved
% OPENAI_CLAIM: {"schema_version":1,"source_repo":"https://github.com/openai/math","source_commit":"adc7f1241b42e322a6451854ab7e4b4c146bf78a","release_date":"2026-10-06","match":"full","resolution":"proved","independently_reviewed":false,"second_pass":"agreed","families":[{"family":"039","title":"Nagata’s conjecture and maximal Seshadri constants","subject":"Algebraic and complex geometry","manuscripts":[{"title":"Maximal Multipoint Seshadri Constants in Higher Dimensions","pdf_path":"preprints/Maximal-Multipoint-Seshadri-Constants-in-Higher-Dimensions-October-5-2026/main.pdf","readme_path":"preprints/Maximal-Multipoint-Seshadri-Constants-in-Higher-Dimensions-October-5-2026/README.md","date":"2026-10-05","theorem_ref":"thm:main (sections/introduction.tex:29)","statement_sources":[{"path":"preprints/Maximal-Multipoint-Seshadri-Constants-in-Higher-Dimensions-October-5-2026/build/sections/introduction.tex","line":29,"label":"thm:main"}]}],"lean":{"doc_path":"lean/docs/039.md","comparators":[],"declarations":[],"covers_main_theorem":false},"reasoning_trace":null}]}
\documentclass[11pt]{article}
\usepackage[margin=25mm]{geometry}
\usepackage{fontspec}
\setmainfont{Latin Modern Roman}
\usepackage{amsmath,amssymb,mathtools}
\usepackage{unicode-math}
\setmathfont{Latin Modern Math}
\usepackage{enumitem}
\usepackage{xurl,hyperref,bookmark}
\hypersetup{colorlinks=true,urlcolor=blue}
\setcounter{secnumdepth}{0}
\setlength{\parindent}{0pt}
\setlength{\parskip}{5pt}
\setlength{\emergencystretch}{3em}
\begin{document}
\section*{Maximal multipoint Seshadri constants in higher complex dimension}

\subsection{OpenAI's claim}
In every dimension at least three, the ordinary multipoint Seshadri constant is claimed to reach $(L^n/r)^{1/n}$ at very general points for every sufficiently large $r$. The cited manuscript argues for that result; its relation to this problem is specified below.

\subsection{Mathematical statement of the claimed result}
\paragraph{Mathematical claim 1.} Let $X$ be smooth integral projective over $\mathbb C$, with $n=\dim X\ge3$, and $L$ ample. OpenAI claims an integer $r_0(X,L)$ such that for every $r\ge r_0$ and very general ordered distinct $p_1,\ldots,p_r$,
\[
\varepsilon(X,L;\mathbf p)=\left(\frac{L^n}{r}\right)^{1/n}.
\]
The multipoint Seshadri constant is the supremum of $t\ge0$ for which $\pi^*L-t\sum_iE_i$ is nef on the blow-up at these points. For each $r$, the exceptional tuples lie in a union, indexed by a countable set, of proper algebraic closed subsets with nonempty complement. The threshold depends on $(X,L)$ and is not specified uniformly.

\subsection{Outline of the argument}
Manuscript 1 (family 039, 2026-10-05): Graded exponent supports are reshaped into a volume-preserving simplex compatible with section multiplication. Full density forces every sufficiently deep interior lattice point to occur, allowing torus interpolation to separate the needed jets; very-general tuples inherit all the finite tests.

\subsection{Scope relative to this problem}
Manuscript 1 (family 039, 2026-10-05): The problem selects the mathematical statement summarized here, including its stated hypotheses and parameter range. The cited theorem gives a claimed proof of that assertion in full.

\subsection{Formal verification}
For family 039, the scope page documents the following selections. The selected plane-curve inequality allows unequal multiplicities and repeated components at every allowed point count. A separate statement gives eventual maximality on each polarized complex surface, with a surface-dependent threshold; the higher-dimensional and positive-characteristic companions are not selected here. Manuscript 1 is not among the papers selected by that page, so the family link does not formalize its main conclusion.

This site has not executed the Lean files or checked their dependencies. These descriptions report the scope of the linked declarations, rather than a new verification of the proofs.

\subsection{Status caveat}
OpenAI attributes these results to an internal model. Its release notice cautions that results lacking full formalization can contain mistakes. This site has not independently assessed the proofs; a solved label records the scope of the reported claim.

\subsection{Sources}
\begin{itemize}
\item OpenAI, \emph{Maximal Multipoint Seshadri Constants in Higher Dimensions}, 2026-10-05; citation key \texttt{OAI:Maximal-Multipoint-Seshadri-Constants-in-Higher-Dimensions-October-5-2026}. \href{https://github.com/openai/math/blob/adc7f1241b42e322a6451854ab7e4b4c146bf78a/preprints/Maximal-Multipoint-Seshadri-Constants-in-Higher-Dimensions-October-5-2026/main.pdf}{PDF (GitHub)}; \href{https://github.com/openai/math/raw/adc7f1241b42e322a6451854ab7e4b4c146bf78a/preprints/Maximal-Multipoint-Seshadri-Constants-in-Higher-Dimensions-October-5-2026/main.pdf}{PDF (direct)}; \href{https://github.com/openai/math/blob/adc7f1241b42e322a6451854ab7e4b4c146bf78a/preprints/Maximal-Multipoint-Seshadri-Constants-in-Higher-Dimensions-October-5-2026/README.md}{Manuscript page and citation}; \href{https://github.com/openai/math/blob/main/preprints/Maximal-Multipoint-Seshadri-Constants-in-Higher-Dimensions-October-5-2026/main.pdf}{Latest version}; Statement source: \href{https://github.com/openai/math/blob/adc7f1241b42e322a6451854ab7e4b4c146bf78a/preprints/Maximal-Multipoint-Seshadri-Constants-in-Higher-Dimensions-October-5-2026/build/sections/introduction.tex\#L29}{thm:main}.
\item Family 039: \href{https://github.com/openai/math/blob/adc7f1241b42e322a6451854ab7e4b4c146bf78a/lean/docs/039.md}{Lean scope}.
\item \href{https://github.com/openai/math/blob/adc7f1241b42e322a6451854ab7e4b4c146bf78a/CONTENTS.md}{OpenAI catalogue entry}.
\end{itemize}
\end{document}
